\chapter{PHƯƠNG PHÁP LUẬN}
\label{chap:methodology}

\section{Tổng quan kiến trúc Framework}
\label{sec:architecture-overview}

Framework được đề xuất implement một pipeline end-to-end để tạo ra các mô hình AI cho theo dõi chuyển động trên edge devices. Kiến trúc hệ thống bao gồm năm thành phần chính:

\begin{enumerate}
    \item \textbf{Thu thập dữ liệu:} Ứng dụng mobile để capture dữ liệu IMU sensor
    \item \textbf{Preprocessing:} X thập đường chuẩn hóa, làm sạch và phân đoạn dữ liệu
    \item \textbf{Feature extraction:} Trích xuất orientation-invariant features
    \item \textbf{Model training:} Huấn luyện SVM và Neural Network models
    \item \textbf{Code generation:} Tự động sinh C++ code cho embedded platforms
\end{enumerate}

\section{Nền tảng phần cứng và Thu thập dữ liệu}
\label{sec:hardware-data-collection}

\subsection{Seeed XIAO nRF52840 Platform}
Hệ thống sử dụng Seeed XIAO nRF52840 Sense làm platform thu thập dữ liệu và deployment:
\begin{itemize}
    \item \textbf{Processor:} ARM Cortex-M4F @ 64MHz
    \item \textbf{Memory:} 256KB RAM, 1MB Flash
    \item \textbf{IMU Sensor:} LSM6DS3 6-axis (3-axis accelerometer + 3-axis gyroscope)
    \item \textbf{Communication:} I2C address 0x6A
    \item \textbf{Power:} 5mW average consumption
\end{itemize}

\subsection{Cấu hình IMU Sensor}
LSM6DS3 được cấu hình với các parameters sau:
\begin{itemize}
    \item Accelerometer range: $\pm$2g (độ nhạy phù hợp cho human motion)
    \item Gyroscope range: $\pm$245 deg/s
    \item Sampling rate: 100Hz (đủ để capture human motion < 20Hz)
    \item Output data rate (ODR): 104Hz
    \item Low-pass filter: Enabled để giảm noise tần số cao
\end{itemize}

\subsection{Protocol Thu thập dữ liệu}
Firmware thu thập dữ liệu thực hiện các bước sau:

\begin{enumerate}
    \item Initialize I2C communication và verify IMU connection
    \item Configure accelerometer và gyroscope registers
    \item Đọc raw sensor readings mỗi 10ms (100Hz)
    \item Convert accelerometer values sang m/s² (factor: 9.80665)
    \item Convert gyroscope values sang degrees/second
    \item Output CSV format: $aX, aY, aZ, gX, gY, gZ, Time_{seconds}$
\end{enumerate}

Dữ liệu được streaming qua USB Serial với baudrate 115200 và lưu trực tiếp vào CSV files trên computer.

\section{Pipeline Xử lý Dữ liệu}
\label{sec:preprocessing-pipeline}

\subsection{Sliding Window Segmentation}
Dữ liệu time-series được chia thành các windows với overlapping:

\begin{itemize}
    \item \textbf{Window size:} 50 samples (0.5 giây @ 100Hz)
    \item \textbf{Overlap:} 50\% (25 samples)
    \item \textbf{Step size:} 25 samples (0.25 giây)
\end{itemize}

Window size 50 samples được chọn dựa trên phân tích:
\begin{itemize}
    \item Đủ dài để capture 1 chu kỳ walking cycle ($\sim$0.5s)
    \item Đủ ngắn để detect activity transitions nhanh
    \item Balance giữa temporal resolution và computational cost
\end{itemize}

Overlap 50\% đảm bảo:
\begin{itemize}
    \item Tăng số training samples (data augmentation)
    \item Smooth transitions giữa các windows
    \item Capture được movements xảy ra ở biên windows
\end{itemize}

\subsection{Orientation-Invariant Feature Extraction}
Để giải quyết vấn đề device orientation, system sử dụng magnitude-based features:

\textbf{Bước 1: Tính Magnitude Vectors}
\begin{align}
\text{acc\_mag}[i] &= \sqrt{aX[i]^2 + aY[i]^2 + aZ[i]^2} \\
\text{gyro\_mag}[i] &= \sqrt{gX[i]^2 + gY[i]^2 + gZ[i]^2}
\end{align}

\textbf{Bước 2: Tính Jerk (tốc độ thay đổi gia tốc)}
\begin{equation}
\text{jerk\_mag}[i] = |\text{acc\_mag}[i] - \text{acc\_mag}[i-1]| \times 100
\end{equation}

\textbf{Bước 3: Trích xuất 17 Statistical Features}

Cho mỗi window, tính các features sau:

\textit{Accelerometer magnitude features (6 features):}
\begin{itemize}
    \item $f_0$: acc\_mag\_mean = $\frac{1}{N}\sum_{i=1}^{N} \text{acc\_mag}[i]$
    \item $f_1$: acc\_mag\_std = $\sqrt{\frac{1}{N}\sum_{i=1}^{N}(\text{acc\_mag}[i] - \mu)^2}$
    \item $f_2$: acc\_mag\_rms = $\sqrt{\frac{1}{N}\sum_{i=1}^{N}\text{acc\_mag}[i]^2}$
    \item $f_3$: acc\_mag\_energy = $\sum_{i=1}^{N}\text{acc\_mag}[i]^2$
    \item $f_{12}$: acc\_mag\_min = $\min(\text{acc\_mag})$
    \item $f_{13}$: acc\_mag\_max = $\max(\text{acc\_mag})$
\end{itemize}

\textit{Gyroscope magnitude features (6 features):}
\begin{itemize}
    \item $f_4$ đến $f_7$: gyro\_mag mean, std, rms, energy
    \item $f_{14}$, $f_{15}$: gyro\_mag min, max
\end{itemize}

\textit{Jerk magnitude features (5 features):}
\begin{itemize}
    \item $f_8$ đến $f_{11}$: jerk\_mag mean, std, rms, energy
    \item $f_{16}$: jerk\_mag\_max
\end{itemize}

\textbf{Feature Order Critical:} Thứ tự features phải khớp chính xác giữa training và deployment:
\begin{equation}
\mathbf{F} = [f_0, f_1, ..., f_{11}, f_{12}, f_{13}, f_{14}, f_{15}, f_{16}]
\end{equation}

\subsection{Feature Normalization}
Sau khi extract features, áp dụng StandardScaler normalization:

\begin{equation}
f'_i = \frac{f_i - \mu_i}{\sigma_i}
\end{equation}

Trong đó:
\begin{itemize}
    \item $f_i$: raw feature value
    \item $\mu_i$: mean của feature $i$ trên toàn bộ training set
    \item $\sigma_i$: standard deviation của feature $i$
    \item $f'_i$: normalized feature value
\end{itemize}

\textbf{Lưu ý quan trọng:} Scaler phải được fit trên RAW features (chưa normalize), không được fit trên dữ liệu đã normalize trước đó để tránh double normalization bug.

\section{Machine Learning Models}
\label{sec:ml-models}

\subsection{Support Vector Machine (SVM)}
\subsubsection{Thuật toán}
SVM tìm hyperplane tối ưu phân tách các classes:

\begin{equation}
\mathbf{w}^T\mathbf{x} + b = 0
\end{equation}

Với RBF (Radial Basis Function) kernel:
\begin{equation}
K(\mathbf{x}_i, \mathbf{x}_j) = \exp(-\gamma ||\mathbf{x}_i - \mathbf{x}_j||^2)
\end{equation}

\subsubsection{Hyperparameters}
\begin{itemize}
    \item \textbf{C (regularization):} 1.0 (default) - balance margin và misclassification
    \item \textbf{gamma:} 'scale' = $\frac{1}{n_{features} \times \text{var}(X)}$
    \item \textbf{kernel:} 'rbf' - xử lý tốt non-linear relationships
    \item \textbf{decision\_function\_shape:} 'ovr' (one-vs-rest) cho multi-class
\end{itemize}

\subsubsection{Training Process}
\begin{enumerate}
    \item Load normalized features và labels
    \item Split data: 70\% train, 15\% validation, 15\% test
    \item Fit SVM model với GridSearchCV để tìm best params
    \item Evaluate trên validation set
    \item Final evaluation trên test set
\end{enumerate}

\subsection{Multi-layer Neural Network}
\subsubsection{Kiến trúc}
3-layer feedforward neural network:

\begin{itemize}
    \item \textbf{Input layer:} 17 neurons (17 features)
    \item \textbf{Hidden layer 1:} 100 neurons, ReLU activation
    \item \textbf{Hidden layer 2:} 50 neurons, ReLU activation
    \item \textbf{Output layer:} 5 neurons (5 classes), Linear activation
\end{itemize}

Architecture: 17 $\rightarrow$ 100 $\rightarrow$ 50 $\rightarrow$ 5

\subsubsection{Forward Propagation}

\textbf{Layer 1 (Input $\rightarrow$ Hidden1):}
\begin{align}
\mathbf{z}^{(1)} &= \mathbf{W}^{(1)}\mathbf{x} + \mathbf{b}^{(1)} \\
\mathbf{a}^{(1)} &= \text{ReLU}(\mathbf{z}^{(1)}) = \max(0, \mathbf{z}^{(1)})
\end{align}

\textbf{Layer 2 (Hidden1 $\rightarrow$ Hidden2):}
\begin{align}
\mathbf{z}^{(2)} &= \mathbf{W}^{(2)}\mathbf{a}^{(1)} + \mathbf{b}^{(2)} \\
\mathbf{a}^{(2)} &= \text{ReLU}(\mathbf{z}^{(2)})
\end{align}

\textbf{Layer 3 (Hidden2 $\rightarrow$ Output):}
\begin{equation}
\mathbf{y} = \mathbf{W}^{(3)}\mathbf{a}^{(2)} + \mathbf{b}^{(3)}
\end{equation}

Prediction:
\begin{equation}
\text{class} = \arg\max_i(\mathbf{y}_i)
\end{equation}

\subsubsection{Hyperparameters}
\begin{itemize}
    \item \textbf{Solver:} 'adam' - adaptive learning rate
    \item \textbf{Learning rate:} 0.001 (initial)
    \item \textbf{Max iterations:} 500
    \item \textbf{Batch size:} 'auto' = min(200, n\_samples)
    \item \textbf{Activation:} 'relu' (hidden), linear (output)
    \item \textbf{Alpha (L2 regularization):} 0.0001
\end{itemize}

\section{Automatic Code Generation}
\label{sec:code-generation}

\subsection{Architecture Overview}
Code generator tự động chuyển đổi trained scikit-learn models sang optimized C++ code.

\textbf{Components:}
\begin{itemize}
    \item \textbf{BaseGenerator:} Abstract class với common functionality
    \item \textbf{SVMGenerator:} Generate SVM-specific code
    \item \textbf{NeuralNetworkGenerator:} Generate NN-specific code
    \item \textbf{Template Engine:} Sử dụng string formatting với placeholders
\end{itemize}

\subsection{SVM Code Generation}

\textbf{Bước 1: Extract Model Parameters}
\begin{itemize}
    \item Support vectors: $\mathbf{SV} \in \mathbb{R}^{n_{sv} \times 17}$
    \item Dual coefficients: $\alpha_i$ for each support vector
    \item Intercepts: $b_j$ for each class pair
    \item Kernel parameters: $\gamma$
\end{itemize}

\textbf{Bước 2: Generate Prediction Function}
\begin{verbatim}
int svm_predict(float features[17]) {
    // Step 1: Normalize features
    float scaled[17];
    for(int i=0; i<17; i++)
        scaled[i] = (features[i] - means[i]) / stds[i];
    
    // Step 2: Compute RBF kernel với support vectors
    float decision_values[num_classes];
    for(int c=0; c<num_classes; c++) {
        float sum = 0;
        for(int sv=0; sv<n_support_vectors; sv++) {
            float diff_sq = 0;
            for(int f=0; f<17; f++) {
                float diff = scaled[f] - support_vectors[sv][f];
                diff_sq += diff * diff;
            }
            float kernel = exp(-gamma * diff_sq);
            sum += dual_coef[sv] * kernel;
        }
        decision_values[c] = sum + intercepts[c];
    }
    
    // Step 3: Return class with highest score
    return argmax(decision_values);
}
\end{verbatim}

\subsection{Neural Network Code Generation}

\textbf{Bước 1: Extract Weight Matrices}
\begin{itemize}
    \item $\mathbf{W}^{(1)} \in \mathbb{R}^{17 \times 100}$: input $\rightarrow$ hidden1
    \item $\mathbf{W}^{(2)} \in \mathbb{R}^{100 \times 50}$: hidden1 $\rightarrow$ hidden2
    \item $\mathbf{W}^{(3)} \in \mathbb{R}^{50 \times 5}$: hidden2 $\rightarrow$ output
    \item Biases: $\mathbf{b}^{(1)}, \mathbf{b}^{(2)}, \mathbf{b}^{(3)}$
\end{itemize}

\textbf{Bước 2: Generate C++ Arrays}
\begin{verbatim}
const float input_weights[17][100] = { ... };
const float hidden_biases[100] = { ... };
const float hidden2_weights[100][50] = { ... };
const float hidden2_biases[50] = { ... };
const float final_weights[50][5] = { ... };
const float final_biases[5] = { ... };
\end{verbatim}

\textbf{Bước 3: Generate Inference Function}
\begin{verbatim}
int nn_predict(float features[17]) {
    // Normalize
    float scaled[17];
    normalize(features, scaled);
    
    // Layer 1
    float hidden1[100];
    for(int h=0; h<100; h++) {
        float sum = hidden_biases[h];
        for(int i=0; i<17; i++)
            sum += scaled[i] * input_weights[i][h];
        hidden1[h] = relu(sum);
    }
    
    // Layer 2
    float hidden2[50];
    for(int h=0; h<50; h++) {
        float sum = hidden2_biases[h];
        for(int i=0; i<100; i++)
            sum += hidden1[i] * hidden2_weights[i][h];
        hidden2[h] = relu(sum);
    }
    
    // Output layer
    float output[5];
    for(int o=0; o<5; o++) {
        float sum = final_biases[o];
        for(int h=0; h<50; h++)
            sum += hidden2[h] * final_weights[h][o];
        output[o] = sum;
    }
    
    return argmax(output);
}
\end{verbatim}

\subsection{Memory Optimization Techniques}

\textbf{1. Fixed-point Arithmetic (Optional):}
Chuyển float32 sang int16:
\begin{equation}
\text{int16\_value} = \text{round}(\text{float\_value} \times 2^{14})
\end{equation}

Giảm memory 50\% nhưng cần verify accuracy impact.

\textbf{2. Weight Pruning:}
Loại bỏ weights < threshold (e.g., $|w| < 0.01$):
\begin{itemize}
    \item Giảm computation
    \item Tăng sparsity cho compression
    \item Trade-off: có thể giảm accuracy
\end{itemize}

\textbf{3. Const Arrays in Flash:}
Sử dụng \texttt{const} keyword để lưu weights vào Flash thay vì RAM:
\begin{verbatim}
const float PROGMEM weights[100][50] = { ... };
\end{verbatim}

\section{Deployment Workflow}
\label{sec:deployment-workflow}

\subsection{Complete Pipeline}
\begin{enumerate}
    \item \textbf{Data Collection:} Thu thập IMU data với XIAO, lưu CSV
    \item \textbf{Preprocessing:} Load CSV, sliding window, extract features
    \item \textbf{Training:} Train SVM/NN, optimize hyperparameters
    \item \textbf{Validation:} Evaluate accuracy, confusion matrix
    \item \textbf{Code Generation:} Generate C++ .h và .cpp files
    \item \textbf{Integration:} Include vào Arduino project
    \item \textbf{Compilation:} Build với Arduino IDE
    \item \textbf{Upload:} Flash lên XIAO board
    \item \textbf{Testing:} Verify real-time predictions
\end{enumerate}

\subsection{Generated Code Structure}
\begin{verbatim}
har_model_xiao_f17_c5/
├── har_model_xiao_f17_c5.h      # Header với declarations
├── har_model_xiao_f17_c5.cpp    # Implementation
└── README.md                     # Usage instructions
\end{verbatim}

\textbf{Header file contains:}
\begin{itemize}
    \item Model constants (NUM\_FEATURES, NUM\_CLASSES)
    \item Function declarations (har\_predict, har\_extract\_features)
    \item Activity name mappings
\end{itemize}

\textbf{Implementation file contains:}
\begin{itemize}
    \item Feature scaling parameters (means, stds)
    \item Model weights (SVM support vectors hoặc NN weights)
    \item Feature extraction logic
    \item Prediction function
\end{itemize}

\section{Tóm tắt}
Chương này đã trình bày chi tiết methodology của framework, bao gồm:
\begin{itemize}
    \item Hardware platform và data collection protocol
    \item Preprocessing pipeline với orientation-invariant features
    \item Machine learning algorithms (SVM, Neural Network)
    \item Automatic code generation cho embedded deployment
    \item Optimization techniques để giảm memory footprint
\end{itemize}

Chương tiếp theo sẽ present kết quả thử nghiệm và đánh giá hiệu năng của framework.
